\section{Appendix D. Spectral Harmonics and the Cymatic Analogy in Intrinsic Response Theory}

This appendix formalizes the ``cymatics'' analogy used throughout the paper by identifying a single spectral object---the eigen-spectrum of an intrinsic response operator---that plays the same mathematical role as Laplacian/biharmonic eigenmodes in Chladni plate experiments. In brief: eigenfunctions are the response-sector mode shapes, eigenvalues set the characteristic spatial scales, and nodal sets correspond to compression/rarefaction boundaries. Although the underlying response dynamics may be written covariantly on spacetime, the galaxy-rotation regime probed by SPARC is effectively stationary, so the empirically constrained ``harmonics'' are spatial overtones (additional radial nodes) rather than temporal frequency multiples.

\subsection{Purpose and scope}

The goal is to tie the qualitative ``pattern'' language to a concrete operator-theoretic object that can be constrained by the rotation-curve phenomenology.

\subsection{Unified operator statement (“DNA” of intrinsic response)}

Let \textbf{(see PDF; response-potential symbol not recovered)} denote the response-sector potential (or any equivalent scalar response field used to encode the inferred metric perturbation in the weak-field limit). Define the intrinsic response through a driven, self-adjoint elliptic operator on a spatial slice:

\textbf{(see PDF; driven elliptic operator equation not recovered from DOCX extraction)}

where \textbf{(see PDF; constitutive driving functional not recovered)} is the constitutive driving functional built from the baryonic sector, and the operator is taken in a minimal Sturm--Liouville form:

\textbf{(see PDF; Sturm--Liouville operator form not recovered from DOCX extraction)}

with \textbf{(see PDF; propagation coefficient not recovered)} controlling propagation (``permittivity/elasticity'') and $\beta(x)\ge 0$ encoding activation/screening (``mass'' or quenching). The intrinsic ``pattern DNA'' is the eigen-spectrum

\textbf{(see PDF; eigenproblem not recovered from DOCX extraction)}

yielding the modal response

\textbf{(see PDF; modal expansion not recovered from DOCX extraction)}

This is the direct analog of Chladni/cymatics: the eigenfunctions are mode shapes; eigenvalues set intrinsic spatial scales; and source overlap coefficients determine which modes are visible in a given galaxy.

\subsection{Where “harmonics” live in a unified spacetime theory}

In a unified spacetime theory, the fundamental response dynamics can be written covariantly (e.g., for a scalar template, a hyperbolic operator such as $\Box$ acting on a field). In stationary or quasi-stationary systems, variable separation reduces the covariant problem to a spatial eigenproblem: \textbf{(see PDF; separation-of-variables equations not recovered)}. SPARC rotation curves constrain the near-static regime in which the response has equilibrated, leaving the spatial spectral content as the observable imprint.

\subsection{Observable proxy: node count and Δ-sign flips in the SPARC window}

SPARC provides baryonic mass models and high-quality rotation curves for 175 disk galaxies. Define \textbf{(see PDF; residual definition and symbol not recovered)} operationally as the signed residual between the observed and baryonic expectations in the rotation curve (e.g., \textbf{see PDF; example residual expression}). In an oscillatory response sector, this residual (and therefore the inferred extra acceleration) can alternate in sign across radius, producing compression and rarefaction lobes. In the operator language above, each radial sign change corresponds to crossing a nodal annulus/shell of the net excited response pattern. Hence, within the finite SPARC radial sampling window, the number of sign flips provides a model-agnostic proxy for the degree of higher-(harmonic/overtone) participation: zero flips indicates a fundamental-dominated response; multiple flips indicate increasing harmonic content (higher node count) selected by the driving and by activation/quenching boundaries.

\subsection{Consequences for lensing and stability constraints}

Because the response tensor is inferred effectively from kinematics, sign-alternating lobes should be treated as an empirical prediction of the effective sector rather than as an immediate pathology. In particular, annular rarefaction phases imply a ring-like suppression of weak-lensing observables (e.g., \textbf{see PDF; observable symbols not recovered}) in the same radial ranges, providing a falsifiable discriminant using upcoming wide-field surveys. At the fundamental level, the minimal non-negotiable stability line is to keep the underlying response-field EFT free of ghosts/gradient instabilities; for broad $k$-essence classes this is ensured by requiring a positive kinetic coefficient \textbf{(see PDF; positivity/stability condition not recovered)}, which enforces the null energy condition for the fundamental field stress--energy even if the effective, inverse-constructed response density changes sign locally. For context, relativistic MOND completions (e.g., TeVeS) also introduce extra fields that alter lensing, but do not generically predict phase-correlated annular suppression tied to sign flips.
